\section{Discussion}\label{sec:discussion}

\para{Implications}
For strong tree ensembles with a tuned threshold, \emph{which} oversampler is applied, and \emph{where}, matters
little (gaps G1--G2). The exception is aggressive cleaning: SMOTE-ENN, and to a lesser degree ADASYN, reliably hurt.
Two mechanisms explain this. First, a strong learner combined with a tuned threshold already makes the
precision--recall trade-off that resampling targets~\cite{elor2022smote,vandengoorbergh2022harm}. Second, once
budgets are matched and leakage removed, most resamplers produce near-identical training sets. Gap G3 fares better:
multiclass problems and datasets with a few dozen minority examples show the largest gains, and resampling research
should focus there.

\para{Five evaluation pitfalls}
Our first pipeline contained five errors, each of which would have produced a spurious improvement:
\begin{enumerate}
  \item \textbf{Unequal budgets.} Local balancing added far less synthetic data than global SMOTE (a 1:200 versus
        1:1 final ratio on creditcard), confounding resampler choice with resampling amount.
  \item \textbf{Scale-cancelled weights.} Confidence weights fed into a meta-learner are a no-op.
  \item \textbf{Rank-forced rules.} Percentile rules with one or two clusters decide by rank.
  \item \textbf{Fixed 0.5 thresholds.} These make F1, G-mean and MCC mirror the oversampling ratio.
  \item \textbf{Duplicates and pre-split resampling.} Both place near-copies of test points in the training data.
\end{enumerate}
Controlling for all five is, in our view, a minimum standard for resampling papers.

\para{Threats to validity}
\begin{itemize}
  \item \textbf{Fixed hyperparameters.} Hyperparameters and rule thresholds were fixed a priori, not tuned per
        dataset. Tuning would affect all variants alike.
  \item \textbf{Hardware non-determinism.} It exists (GPU versus CPU XGBoost), so the study ran on one machine from
        one code version.
  \item \textbf{Small validation splits.} The 15\% split is small for tiny minority classes.
  \item \textbf{External validity.} The benchmarks form a single collection. Deduplication limits comparability with
        prior work.
  \item \textbf{Learner scope.} The conclusions concern strong tree ensembles; weaker learners may benefit more from
        resampling~\cite{vanhulse2007experimental}.
  \item \textbf{Statistical power.} Per-dataset power is limited by 15 correlated folds, and the minority-size trend
        does not survive multiplicity correction.
\end{itemize}

\section{Future Work}\label{sec:future}
\begin{itemize}
  \item \textbf{Other learners.} Repeat the study with logistic regression, $k$-NN, SVMs and deep tabular models, to
        map the learner--resampler interaction.
  \item \textbf{The small-sample regime.} Test the minority-size trend in a pre-registered study on datasets with
        fewer than 100 minority examples, including multiclass ones, and add generative oversamplers.
  \item \textbf{Meta-learned selection.} Replace rules and noisy inner validation with a meta-model that predicts the
        benefit of resampling from dataset characteristics.
  \item \textbf{Combination.} Replace global confidence weights with dynamic ensemble selection~\cite{cruz2018dynamic}.
  \item \textbf{Alternatives to resampling.} Compare resampling with cost-sensitive learning, and evaluate
        calibration~\cite{vandengoorbergh2022harm}.
  \item \textbf{Analysis and deployment.} Use hierarchical Bayesian models across datasets~\cite{benavoli2017time},
        and add drift monitoring and energy accounting~\cite{plan_p10_stacking} to the served models.
\end{itemize}

\section{Conclusion}\label{sec:conclusion}
We proposed \ars{}, which selects a resampler per minority cluster from overlap and density scores and combines three
tree ensembles with confidence weights, and we evaluated it in a budget-matched, leakage-free study on 28 datasets.

\ars{} reliably beats aggressive cleaning (SMOTE-ENN). It is practically equivalent to plain SMOTE and to its own
ablations, and it cannot be separated from training without resampling. Each null result has a concrete cause:
near-uniform weights, cleaning that removes nothing after balancing, and selection that rarely varies within a
dataset. A validated-selection pilot confirms that the headroom is small, and resampling mainly pays off when minority
examples are scarce.

We hope the method, the pipeline and the five pitfalls help future work target the settings where resampling matters,
and measure its benefit fairly.

\para{Availability}
Code, configuration, per-fold results and analysis scripts are in the project repository. All datasets are public:
Kaggle~\cite{dalpozzolo2015calibrating}, imbalanced-learn~\cite{lemaitre2017imbalanced} and UCI.
